%========================
% Chapter 3 — Methodology (Proposal)
%========================

% \chapter{Methodology} % uncomment if your class uses \chapter

This chapter details the methodological framework of \textit{Sonic Speculocultural Technopoiesis} (ST) and describes the research orientation, cohort design, procedures, data, analysis, rigor, ethics, and deliverables. ST integrates critical technical practice, research-creation, autoethnography, and small-cohort collaboration to operationalize culturally grounded tool modification in audio interaction.

\subsubsection{Objectives}
\begin{itemize}[nosep]
	\item Define ST as an actionable method joining critique and implementation.
	\item Specify cohort composition, procedures, instruments, and data.
	\item Detail a qualitative analysis plan that links cultural rationale to computational change.
	\item Establish evaluation criteria, rigor strategies, and ethical safeguards.
	\item Enumerate deliverables and a feasible timeline.
\end{itemize}

\subsubsection{Intended Sections/Subsections}
\begin{itemize}[nosep]
	\item Research Orientation and Design Logic
		\begin{itemize}[nosep]
			\item research-creation (artistic making as knowledge production)
			\item autoethnography (researcher practice as data)
			\item collaborative small-cohort design (co-creation with artist-technologists)
		\end{itemize}
	\item Framework Mechanics (ST Cycle and Domains)
		\begin{itemize}[nosep]
			\item \textbf{Unsettling}: elicit friction between cultural knowledge and tool defaults; identify leverage points.
			\item \textbf{Remixing}: modify representations, parameters, and algorithms to encode Black sonic logics.
			\item \textbf{Embodying}: assess cultural authenticity and experiential fit through practice and performance.
			\item \textbf{Transmitting}: document methods, name patterns, and package code for reuse.
			\item Design decisions are additionally guided by four cross-cutting territories:
			\textit{Fugitive Speculation}, \textit{Material Memory}, \textit{Sensory Worlding}, and \textit{Intimate Witnessing}. These act as evaluative lenses in every iteration.
		\end{itemize}
	\item Participants and Recruitment
	\item Sites, Materials, and Instruments
	\item Procedures
	\item Data Sources and Analysis Plan
	\item Evaluation Criteria
	\item Rigor and Validity
	\item Ethics, Risk, and Data Security
	\item Deliverables and Timeline	
\end{itemize}

\begin{comment}
%------------------------------------------------
\subsubsection{Research Orientation and Design Logic}
The project adopts a practice-based, qualitative orientation in which making is a primary site of inquiry. Tool prototypes and modifications are treated as research instruments: they materialize theoretical commitments, reveal hidden assumptions, and enable situated evaluation through practice. The study combines:
\begin{itemize}
	\item critical technical practice (theory-making looped with implementation),
	\item research-creation (artistic making as knowledge production),
	\item autoethnography (researcher practice as data),
	\item collaborative small-cohort design (co-creation with artist-technologists).
\end{itemize}

\subsubsection{Framework Mechanics: ST Cycle and Territories}
ST proceeds through four movements that repeat iteratively:
\begin{itemize}
	\item \textbf{Unsettling}: elicit friction between cultural knowledge and tool defaults; identify leverage points.
	\item \textbf{Remixing}: modify representations, parameters, and algorithms to encode Black sonic logics.
	\item \textbf{Embodying}: assess cultural authenticity and experiential fit through practice and performance.
	\item \textbf{Transmitting}: document methods, name patterns, and package code for reuse.
\end{itemize}
Design decisions are additionally guided by four cross-cutting territories:
\textit{Fugitive Speculation}, \textit{Material Memory}, \textit{Sensory Worlding}, and \textit{Intimate Witnessing}.
These act as evaluative lenses in every iteration.

\subsubsection{Participants and Recruitment}
\textbf{Cohort composition.} A small cohort of Black artist-technologists (target 4-8) with experience in sound, performance, or interactive media. \\
\textbf{Inclusion.} Demonstrated practice with Black sonic traditions and interest in tool co-design. \\
\textbf{Recruitment.} Purposeful sampling via professional networks, prior collaborators, and open calls. \\
\textbf{Compensation.} Stipends or honoraria for sessions and follow-up interviews. \\
\textbf{Consent.} Written informed consent with clear options for anonymity or named credit.

\subsubsection{Sites, Materials, and Instruments}
\textbf{Sites.} Studio, lab, or rehearsal space with quiet recording conditions; remote sessions as needed. \\
\textbf{Materials.} Audio workstation and prototyping tools (e.g., DAW, Max/MSP or equivalent), optional game engine for interaction prototyping, controllers/sensors as needed. \\
\textbf{Instruments.} Screen/audio capture software; microphones/recorder; version control repositories for code and patch files; session templates for prompts and reflection.

\subsubsection{Procedures}
Each cohort engages a structured series of sessions that instantiate the ST cycle.

\subsubsection{Session 0: Baseline and Unsettling Interview}
\begin{itemize}
	\item Warm-up demonstration of current tools and practices.
	\item Prompted reflection on friction with defaults (time bases, representation, control, authorship).
	\item Identification and ranking of leverage points to address first.
\end{itemize}

\subsubsection{Session 1-2: Remixing Sprints}
\begin{itemize}
	\item Co-design and implement one or more modifications (representation, interaction, or algorithm).
	\item Track decisions in a changelog; commit code or patches to a shared repository.
	\item Short playtests at the end of each sprint; capture screen, audio, and commentary.
\end{itemize}

\subsubsection{Session 3: Embodying Test}
\begin{itemize}
	\item Extended tryout in practice conditions (improvisation, movement, performance routine).
	\item Debrief focused on cultural authenticity, affordances, constraints, and emergent techniques.
\end{itemize}

\subsubsection{Session 4: Transmitting and Naming}
\begin{itemize}
	\item Consolidate modifications into named patterns and procedures.
	\item Produce documentation: short pattern writeups, parameter schemas, and annotated examples.
\end{itemize}

\subsubsection{Data Sources}
\begin{itemize}
	\item \textbf{Process data}: design notes, commits, parameter diffs, prototypes and patch files.
	\item \textbf{Performance data}: audio recordings, screen captures, MIDI or control logs.
	\item \textbf{Dialogic data}: session transcripts, interviews, written reflections.
	\item \textbf{Researcher data}: autoethnographic fieldnotes, memos, decision logs.
	\item \textbf{Artifacts for transmission}: pattern library entries, code snippets, diagrams, instructions.
\end{itemize}

\subsubsection{Analysis Plan}
\textbf{Within-case analysis.} For each cohort case, build a timeline of the ST cycle; code transcripts and notes for ST movements and territories; map modifications to the frictions they address; analyze sonic/interaction outcomes. \\
\textbf{Cross-case analysis.} Compare leverage points, recurring patterns, and convergent modifications; derive higher-level pattern language and design principles. \\
\textbf{Tech analysis.} Inspect version diffs, parameter changes, and algorithmic structures to link cultural rationale to computational implementation. \\
\textbf{Codebook.} Initial codes derived from ST (Unsettling, Remixing, Embodying, Transmitting; four territories), plus inductive subcodes emerging from data. Codebook iteratively refined with collaborator feedback.

\subsubsection{Evaluation Criteria}
\begin{itemize}
	\item \textbf{Cultural authenticity}: does the modification enable expression aligned with the cited sonic practice.
	\item \textbf{Expressive affordance}: range and nuance unlocked versus baseline tooling.
	\item \textbf{Friction reduction}: evidence that a specific default was surfaced and materially re-authored.
	\item \textbf{Transmissibility}: clarity and portability of the resulting pattern, code, or procedure.
	\item \textbf{Situated validity}: participant satisfaction, willingness to adopt in practice, and qualitative reports of change.
\end{itemize}

\subsubsection{Rigor and Validity}
\begin{itemize}
	\item \textbf{Triangulation}: converge process, performance, dialogic, and researcher data.
	\item \textbf{Member checks}: participants review transcripts, codes, and pattern writeups.
	\item \textbf{Audit trail}: repositories, changelogs, and memoing document decisions and iterations.
	\item \textbf{Reflexivity}: positionality statements and structured prompts in fieldnotes.
	\item \textbf{Thick description}: richly contextualize sessions, setups, and cultural rationales.
\end{itemize}

\subsubsection{Ethics, Risk, and Data Security}
\textbf{Risks.} Minimal; potential discomfort in critique of existing tools or cultural misrepresentation. \\
\textbf{Mitigation.} Co-authorship of descriptions; opt-in for attribution; right to pause or withdraw at any time. \\
\textbf{Privacy.} Pseudonyms by default; named credit by written request. Audio and video stored on encrypted drives; access limited to research team. \\
\textbf{Data banking.} Research data and artifacts may be banked for future studies with consent. Storage location, retention period, and access controls are specified in the IRB. \\
\textbf{Release procedures.} Requests for access require PI approval per IRB; only de-identified datasets and redacted artifacts are released; participants can decline inclusion.

\subsubsection{Limitations and Delimitations}
Small cohort size and practice-based scope limit generalizability; results are intentionally situated. Technical constraints of platforms may bound what can be modified. The study focuses on audio interaction rather than full production pipelines.

\subsubsection{Deliverables}
\begin{itemize}
	\item Prototypes and code modules implementing culturally grounded interaction or time bases.
	\item A pattern library documenting named modifications, parameter schemas, and example use.
	\item Case study reports with analysis and transcripts.
	\item A methodological guide for applying ST in other contexts.
\end{itemize}

\subsubsection{Timeline}
\begin{itemize}
	\item Phase 1: Recruitment, Session 0 baselines, friction mapping.
	\item Phase 2: Remixing sprints and interim embodying tests.
	\item Phase 3: Final embodying sessions and transmitting documentation.
	\item Phase 4: Analysis, cross-case synthesis, and consolidation of patterns.
\end{itemize}



% --- Chapter 3 Key References (keyword-based filter) ---
\begin{refsection}[references.bib] % ensure your ch.3 items have keyword=ch3
	\nocite{*}
	\printbibliography[keyword=ch3,heading=subbibliography,title={Key References}]
\end{refsection}
\end{comment}